\documentclass[12pt,a4paper]{article}
\usepackage[utf8]{inputenc}
\usepackage[T1]{fontenc}
\usepackage{amsmath,amsfonts,amssymb,amsthm}
\usepackage{geometry}
\usepackage{fancyhdr}
\usepackage{enumerate}
\usepackage{hyperref}
\usepackage{xcolor}

\geometry{margin=2.5cm}
\pagestyle{fancy}
\fancyhf{}
\fancyhead[C]{GPT-5.4 Solutions -- Petropolitan 2026}
\fancyfoot[C]{\thepage}

\title{GPT-5.4 Solutions\\Petropolitan Science (Re)Search 2026}
\author{Generated by imc\_math\_llm\_eval.py}
\date{\today}

\begin{document}

\maketitle
\tableofcontents
\newpage

\section{Problem 1, Petropolitan 2026}
\subsection{Variant 1 (33.0s)}
We seek all pairs \((k,a)\in \mathbb N\times \mathbb R\) for which there exists a continuous \(f:[0,1]\to[0,1]\) satisfying
\[
\sum_{i=0}^{k-1} f_i(x)=a \qquad (\forall x\in[0,1]),
\]
where \(f_0(x)=x\) and \(f_n=f\circ f_{n-1}\).

We will prove that this is possible \textbf{if and only if}
\[
a=\frac{k-1}{2}.
\]

Moreover, for every \(k\), such a function exists; for example \(f(x)=1-x\) works when \(k\) is even, and \(f(x)=\tfrac12\) works when \(k\) is odd.


\medskip\hrule\medskip


\subsection*{1. Reformulation}

Define
\[
S(x):=\sum_{i=0}^{k-1} f_i(x).
\]
The condition is that \(S(x)\equiv a\) is constant on \([0,1]\).

We must determine for which \(a\) this can happen.


\medskip\hrule\medskip


\subsection*{2. Immediate examples}

Before proving necessity, let us note two constructions.

\subsubsection*{Case 1: \(k\) even}
Take
\[
f(x)=1-x.
\]
Then
\[
f_0(x)=x,\quad f_1(x)=1-x,\quad f_2(x)=x,\quad f_3(x)=1-x,\dots
\]
Hence the sum of any two consecutive iterates is \(x+(1-x)=1\). Since \(k\) is even,
\[
S(x)=\frac{k}{2}\cdot 1=\frac{k}{2}.
\]
So in this case \(a=\frac{k}{2}\).

But note that when \(k\) is even,
\[
\frac{k-1}{2}\neq \frac{k}{2},
\]
so this seems contradictory to the claimed answer. Therefore we must be careful: the problem states the sum runs from \(f_0\) to \(f_{k-1}\), so for \(k=2\), indeed \(x+f(x)=1\), giving \(a=1\). That suggests the true answer is not yet identified. So we continue carefully.

\subsubsection*{Case 2: \(k\) odd}
Take
\[
f(x)=\frac12.
\]
Then
\[
f_0(x)=x,\qquad f_i(x)=\frac12 \ \ (i\ge1),
\]
so
\[
S(x)=x+\frac{k-1}{2},
\]
which is not constant. So this does \textbf{not} work.

So we need a real derivation.


\medskip\hrule\medskip


\subsection*{3. A key identity}

Since \(S(x)=a\) for all \(x\), substituting \(f(x)\) for \(x\) gives
\[
\sum_{i=0}^{k-1} f_i(f(x))=a.
\]
But \(f_i(f(x))=f_{i+1}(x)\), so
\[
\sum_{i=1}^{k} f_i(x)=a.
\]
Subtracting from
\[
\sum_{i=0}^{k-1} f_i(x)=a
\]
yields
\[
f_k(x)=f_0(x)=x \qquad (\forall x\in[0,1]).
\]
Thus
\[
f^k=\operatorname{id}_{[0,1]}.
\]

So any admissible \(f\) is a continuous self-map of \([0,1]\) whose \(k\)-th iterate is the identity.


\medskip\hrule\medskip


\subsection*{4. Continuous interval maps with a finite iterate equal to the identity}

We now show that such an \(f\) must satisfy
\[
f^2=\operatorname{id}.
\]
In fact, more strongly, \(f\) is either the identity or a decreasing involution.

\subsubsection*{Lemma}
If \(f:[0,1]\to[0,1]\) is continuous and \(f^k=\operatorname{id}\) for some \(k\ge1\), then \(f\) is injective.

\textbf{Proof.}
If \(f(u)=f(v)\), then applying \(f^{k-1}\) to both sides gives
\[
u=f^k(u)=f^{k-1}(f(u))=f^{k-1}(f(v))=f^k(v)=v.
\]
So \(u=v\). $\qed$

A continuous injective map on an interval is strictly monotone. Hence \(f\) is either strictly increasing or strictly decreasing.


\medskip\hrule\medskip


\subsubsection*{Lemma}
If \(f\) is strictly increasing and \(f^k=\operatorname{id}\), then \(f=\operatorname{id}\).

\textbf{Proof.}
Suppose there exists \(x\) with \(f(x)>x\). Since \(f\) is increasing,
\[
x<f(x)<f^2(x)<\cdots<f^k(x)=x,
\]
impossible.

Similarly, if \(f(x)<x\), then
\[
x>f(x)>f^2(x)>\cdots>f^k(x)=x,
\]
impossible.

Hence \(f(x)=x\) for all \(x\). $\qed$


\medskip\hrule\medskip


\subsubsection*{Lemma}
If \(f\) is strictly decreasing and \(f^k=\operatorname{id}\), then \(f^2=\operatorname{id}\).

\textbf{Proof.}
Since \(f\) is decreasing, \(f^2\) is increasing. Also
\[
(f^2)^m = f^{2m}.
\]
Because \(f^k=\operatorname{id}\), some positive iterate of \(f^2\) is also the identity (for example \(f^{2k}=\operatorname{id}\)). Thus \(f^2\) is an increasing continuous map with a positive iterate equal to the identity. By the previous lemma,
\[
f^2=\operatorname{id}.
\]
$\qed$

Therefore every admissible \(f\) is either

1. \(f(x)=x\), or
2. a decreasing involution: \(f^2(x)=x\).


\medskip\hrule\medskip


\subsection*{5. Compute the constant sum in each case}

We now use the original equation.

\subsubsection*{Case A: \(f(x)=x\)}

Then every iterate is \(x\), so
\[
S(x)=kx.
\]
This is constant only if \(k=0\), impossible in \(\mathbb N\), so this case gives no solutions.

Thus \(f\) must be decreasing and satisfy \(f^2=\operatorname{id}\).


\medskip\hrule\medskip


\subsubsection*{Case B: \(f\) decreasing and \(f^2=\operatorname{id}\)}

Then the iterates alternate:
\[
f_0(x)=x,\quad f_1(x)=f(x),\quad f_2(x)=x,\quad f_3(x)=f(x),\dots
\]

Hence
\[
S(x)=
\begin{cases}
\frac{k}{2}\bigl(x+f(x)\bigr), & k \text{ even},\\[4pt]
\frac{k-1}{2}\bigl(x+f(x)\bigr)+x, & k \text{ odd}.
\end{cases}
\]

We need this to be constant.


\medskip\hrule\medskip


\subsection*{6. Even \(k\)}

If \(k=2m\), then
\[
S(x)=m(x+f(x)).
\]
So \(x+f(x)\) must be constant:
\[
x+f(x)=c.
\]
Since \(f:[0,1]\to[0,1]\), evaluating at \(x=0,1\) and using that \(f\) is decreasing involution forces \(f(x)=1-x\), hence \(c=1\). Indeed from \(f(x)=c-x\) and \(f([0,1])\subseteq[0,1]\), we must have \(c=1\).

Thus
\[
a=S(x)=m=\frac{k}{2}.
\]

And this is attained by \(f(x)=1-x\).

So for even \(k\), the only possible value is
\[
a=\frac{k}{2}.
\]


\medskip\hrule\medskip


\subsection*{7. Odd \(k\)}

Let \(k=2m+1\). Then
\[
S(x)=m(x+f(x))+x.
\]
If this is constant, then
\[
m(x+f(x))+x=a.
\]
Applying \(f\) instead of \(x\), and using \(f^2(x)=x\), we get
\[
m(f(x)+x)+f(x)=a.
\]
Subtracting the two equations:
\[
x-f(x)=0.
\]
So \(f(x)=x\) for all \(x\), contradicting the decreasing case.

Therefore there are \textbf{no solutions for odd \(k\)}.

(Indeed the identity case also failed.)


\medskip\hrule\medskip


\subsection*{8. Conclusion}

The only possible \(k\) are even, and then necessarily
\[
a=\frac{k}{2}.
\]
Conversely, for every even \(k\), \(f(x)=1-x\) works.

So the complete set of pairs is
\[
\boxed{\{(k,a)\in\mathbb N\times\mathbb R:\ k \text{ is even and } a=\tfrac{k}{2}\}.}
\]


\medskip\hrule\medskip


\textbf{Final answer.}
\[
\boxed{\text{There exists such a continuous }f:[0,1]\to[0,1]\text{ iff }k\text{ is even and }a=\frac{k}{2}.}
\]
One example is \(f(x)=1-x\), for which
\[
x+(1-x)+x+(1-x)+\cdots = \frac{k}{2}.
\]
\newpage
\subsection{Variant 2 (28.6s)}
We seek all pairs \((k,a)\in \mathbb N\times \mathbb R\) such that there exists a continuous \(f:[0,1]\to[0,1]\) satisfying
\[
x+f(x)+f_2(x)+\cdots+f_{k-1}(x)=a \qquad (\forall x\in[0,1]),
\]
where \(f_0(x)=x\) and \(f_n=f\circ f_{n-1}\).

We will determine exactly all such pairs.


\medskip\hrule\medskip


\subsection*{1. Reformulation}

Define
\[
S(x):=\sum_{i=0}^{k-1} f_i(x).
\]
The condition is
\[
S(x)=a \qquad (\forall x\in[0,1]).
\]
Thus \(S\) is constant.

We must find all \(k\in\mathbb N\) and \(a\in\mathbb R\) for which this is possible.


\medskip\hrule\medskip


\subsection*{2. First easy observations}

Since each \(f_i(x)\in[0,1]\), we have
\[
0\le S(x)\le k,
\]
so necessarily
\[
a\in[0,k].
\]

Also, evaluating at \(x=0\) and \(x=1\),
\[
a=f(0)+f_2(0)+\cdots+f_{k-1}(0),
\]
\[
a=1+f(1)+f_2(1)+\cdots+f_{k-1}(1).
\]
Subtracting gives
\[
1+\sum_{i=1}^{k-1}\bigl(f_i(1)-f_i(0)\bigr)=0.
\]
This already suggests strong rigidity.


\medskip\hrule\medskip


\subsection*{3. A key identity}

From
\[
x+f(x)+\cdots+f_{k-1}(x)=a,
\]
apply it to \(f(x)\):
\[
f(x)+f_2(x)+\cdots+f_k(x)=a.
\]
Subtracting,
\[
f_k(x)=x \qquad (\forall x\in[0,1]).
\]
So \(f^k=\mathrm{id}\).

This is the central structural consequence.


\medskip\hrule\medskip


\subsection*{4. Continuous self-maps of an interval with finite order}

We now prove:

\subsubsection*{Lemma}
If \(f:[0,1]\to[0,1]\) is continuous and \(f^k=\mathrm{id}\) for some \(k\ge1\), then \(f=\mathrm{id}\).

\subsubsection*{Proof}
Since \(f^k=\mathrm{id}\), \(f\) is bijective. A continuous bijection from \([0,1]\) to itself is strictly monotone.

So \(f\) is either strictly increasing or strictly decreasing.

- If \(f\) is strictly increasing and \(f\neq \mathrm{id}\), then there exists \(x_0\) with \(f(x_0)\neq x_0\).
  - If \(f(x_0)>x_0\), then by monotonicity
    \[
    x_0<f(x_0)<f^2(x_0)<\cdots<f^k(x_0),
    \]
    hence \(x_0<f^k(x_0)\), contradiction since \(f^k(x_0)=x_0\).
  - If \(f(x_0)<x_0\), similarly
    \[
    x_0>f(x_0)>f^2(x_0)>\cdots>f^k(x_0),
    \]
    contradiction.

  Therefore an increasing \(f\) with \(f^k=\mathrm{id}\) must be \(f=\mathrm{id}\).

- If \(f\) is strictly decreasing, then \(f^2\) is strictly increasing. Also
  \[
  (f^2)^k=f^{2k}=(f^k)^2=\mathrm{id}.
  \]
  By the increasing case, \(f^2=\mathrm{id}\).

  But a strictly decreasing involution on \([0,1]\) certainly exists (for example \(x\mapsto 1-x\)); however if \(f^2=\mathrm{id}\), then our original sum condition will force more. We postpone this for a moment.

Actually we can rule out the decreasing case directly using the constant-sum equation: if \(f\) is decreasing, then \(f^2\) is increasing, and since \(f^k=\mathrm{id}\), the parity matters:
- if \(k\) is odd, \(f^k\) is decreasing, impossible because \(\mathrm{id}\) is increasing;
- if \(k\) is even, then \(f^k\) is increasing, possible in principle, but then \(f^2=\mathrm{id}\).

So the only nontrivial possibility from \(f^k=\mathrm{id}\) alone is a decreasing involution when \(k\) is even. We will show next that the sum condition excludes it unless it coincides with the identity case \(k=1\), which it does not.

This completes the structural reduction. $\qed$


\medskip\hrule\medskip


\subsection*{5. Use the sum condition with \(f^k=\mathrm{id}\)}

We now examine the possibilities.

\subsubsection*{Case 1: \(f=\mathrm{id}\)}

Then
\[
f_i(x)=x \quad \text{for all }i,
\]
so
\[
S(x)=kx.
\]
This is constant on \([0,1]\) only if \(k=0\), but \(k\in\mathbb N\) and in the usual convention \(k\ge1\). For \(k\ge1\), \(kx\) is not constant.

Thus no solution arises from \(f=\mathrm{id}\).

So if a solution exists, it would have to come from the decreasing involution possibility.


\medskip\hrule\medskip


\subsubsection*{Case 2: \(f\) is decreasing and \(f^2=\mathrm{id}\)}

Then \(k\) must be even, say \(k=2m\), because \(f^k=\mathrm{id}\).

Now the iterates alternate:
\[
f_{2r}(x)=x,\qquad f_{2r+1}(x)=f(x).
\]
Hence
\[
S(x)=m\,x+m\,f(x)=m\bigl(x+f(x)\bigr).
\]
So the condition \(S(x)=a\) is equivalent to
\[
x+f(x)=c:=\frac a m \qquad (\forall x\in[0,1]).
\]
Thus necessarily
\[
f(x)=c-x.
\]

Since \(f:[0,1]\to[0,1]\), for all \(x\in[0,1]\) we need \(c-x\in[0,1]\). In particular:
- at \(x=0\): \(c\in[0,1]\),
- at \(x=1\): \(c-1\in[0,1]\), i.e. \(c\in[1,2]\).

Therefore \(c=1\), and hence
\[
f(x)=1-x.
\]
Then
\[
S(x)=m(x+1-x)=m.
\]
So
\[
a=m=\frac k2.
\]

Thus for every even \(k\), the function \(f(x)=1-x\) gives a solution with
\[
a=\frac k2.
\]


\medskip\hrule\medskip


\subsection*{6. Show there are no other solutions}

We have already shown any solution must satisfy \(f^k=\mathrm{id}\). For a continuous self-map of \([0,1]\), this implies \(f\) is a homeomorphism, hence monotone.

- If \(f\) is increasing, then \(f=\mathrm{id}\), impossible.
- Hence \(f\) is decreasing. Then \(k\) must be even, and \(f^2=\mathrm{id}\).
- Therefore the iterates alternate, and the constant-sum condition gives \(x+f(x)=\) constant.
- Since \(f:[0,1]\to[0,1]\), this constant must be \(1\), so \(f(x)=1-x\), and \(a=k/2\).

So the above family is the only one.


\medskip\hrule\medskip


\subsection*{7. Conclusion}

The pairs \((k,a)\) for which such a continuous \(f\) exists are exactly those with \(k\) even and
\[
a=\frac k2.
\]

Indeed, for any even \(k\), the function
\[
f(x)=1-x
\]
satisfies
\[
f_{2r}(x)=x,\qquad f_{2r+1}(x)=1-x,
\]
hence
\[
x+f(x)+\cdots+f_{k-1}(x)=\frac k2\,x+\frac k2\,(1-x)=\frac k2.
\]


\medskip\hrule\medskip


\textbf{Final answer.}  
\[
\boxed{\text{Exactly all pairs }(k,a)\text{ with }k\text{ even and }a=\frac{k}{2}.}
\]
A corresponding function is \(\boxed{f(x)=1-x}\).
\newpage
\subsection{Variant 3 (17.7s)}
We seek all pairs \((k,a)\in \mathbb N\times \mathbb R\) for which there exists a continuous \(f:[0,1]\to[0,1]\) satisfying

\[
\sum_{i=0}^{k-1} f_i(x)=a \qquad (\forall x\in[0,1]),
\]
where \(f_0(x)=x\) and \(f_n=f\circ f_{n-1}\).

We will determine all such pairs completely.


\medskip\hrule\medskip


\subsection*{1. Reformulation}

Define
\[
S_k(x):=x+f(x)+f_2(x)+\cdots+f_{k-1}(x).
\]
We are given that
\[
S_k(x)\equiv a \quad \text{on }[0,1].
\]

We must find all \(k\in\mathbb N\) and \(a\in\mathbb R\) for which such a continuous \(f\) exists.

Since each term lies in \([0,1]\), clearly
\[
0\le a\le k.
\]
But we need the exact possibilities.


\medskip\hrule\medskip


\subsection*{2. First easy cases}

\subsubsection*{Case \(k=1\)}

Then the condition is just
\[
f_0(x)=x=a \qquad \forall x\in[0,1],
\]
impossible since \(x\) is not constant.

So:

\[
k\neq 1.
\]


\medskip\hrule\medskip


\subsection*{3. A key identity}

Since
\[
x+f(x)+\cdots+f_{k-1}(x)=a,
\]
replacing \(x\) by \(f(x)\) gives
\[
f(x)+f_2(x)+\cdots+f_k(x)=a.
\]
Subtracting,
\[
f_k(x)=x \qquad \forall x\in[0,1].
\]

So any such \(f\) must satisfy
\[
f^k=\operatorname{id}_{[0,1]}.
\]

This is very restrictive.


\medskip\hrule\medskip


\subsection*{4. Continuous self-maps of an interval with finite order}

We now prove:

\subsubsection*{Lemma}
If \(f:[0,1]\to[0,1]\) is continuous and \(f^k=\operatorname{id}\) for some \(k\ge 1\), then \(f=\operatorname{id}\).

\paragraph{Proof}
Suppose not. Then there exists \(x_0\in[0,1]\) with \(f(x_0)\ne x_0\).

If \(f(x_0)>x_0\), then by continuity,
\[
f^n(x_0)>f^{n-1}(x_0)>\cdots>f(x_0)>x_0
\]
as long as equality never occurs earlier. In particular the orbit is strictly increasing, so \(f^k(x_0)>x_0\), contradicting \(f^k(x_0)=x_0\).

Similarly, if \(f(x_0)<x_0\), then the orbit is strictly decreasing, hence \(f^k(x_0)<x_0\), again a contradiction.

Therefore \(f(x)=x\) for all \(x\), so \(f=\operatorname{id}\). $\qed$

A slightly more conceptual argument is that a continuous injective map on an interval is monotone, and if \(f^k=\mathrm{id}\), then \(f\) is injective; if increasing, it must be the identity, and if decreasing, then \(f^2=\mathrm{id}\), impossible on \([0,1]\) unless fixed everywhere. But the above direct proof suffices.

Hence from the functional equation we conclude:

\[
f(x)=x \qquad \forall x\in[0,1].
\]


\medskip\hrule\medskip


\subsection*{5. Determine \(a\)}

If \(f(x)=x\), then
\[
f_i(x)=x \quad \text{for all }i,
\]
so
\[
x+f(x)+\cdots+f_{k-1}(x)=kx.
\]
For this to be constant in \(x\), we would need \(kx=a\) for all \(x\in[0,1]\), impossible for every \(k\in\mathbb N\).

Thus no such pair exists.


\medskip\hrule\medskip


\subsection*{6. Check for any overlooked exceptions}

- \(k=1\): impossible.
- \(k\ge 2\): the equation implies \(f^k=\mathrm{id}\), hence \(f=\mathrm{id}\), hence the sum is \(kx\), not constant.

So indeed there are no solutions.


\medskip\hrule\medskip


Final answer.

\[
\boxed{\text{There are no pairs }(k,a)\in\mathbb N\times\mathbb R\text{ for which such a continuous }f\text{ exists.}}
\]
\newpage
\subsection{Final Synthesis (15.1s, total 94.4s)}
Among the drafts, \textbf{Solution 2 is essentially correct}.  
\textbf{Solution 1} has a false start but eventually reaches the right conclusion.  
\textbf{Solution 3} is incorrect: it wrongly claims that a continuous interval map with \(f^k=\mathrm{id}\) must be the identity, overlooking decreasing involutions such as \(f(x)=1-x\).

Here is a concise, corrected complete solution.


\medskip\hrule\medskip


\[
\text{Let }S(x):=\sum_{i=0}^{k-1} f_i(x), \qquad f_0(x)=x,\quad f_{n+1}=f\circ f_n.
\]
We are given that
\[
S(x)=a \qquad \forall x\in[0,1].
\]

We determine all possible pairs \((k,a)\).

\subsection*{1. First consequence: \(f^k=\mathrm{id}\)}

Apply the identity \(S(x)=a\) at the point \(f(x)\). Since \(f_i(f(x))=f_{i+1}(x)\), we get
\[
\sum_{i=1}^{k} f_i(x)=a.
\]
Subtracting from
\[
\sum_{i=0}^{k-1} f_i(x)=a
\]
yields
\[
f_k(x)=x \qquad \forall x\in[0,1].
\]
Hence
\[
f^k=\mathrm{id}_{[0,1]}.
\]

\subsection*{2. Structure of such continuous maps}

From \(f^k=\mathrm{id}\), \(f\) is injective: if \(f(u)=f(v)\), then applying \(f^{k-1}\) gives \(u=v\).  
A continuous injective map on an interval is strictly monotone, so \(f\) is either strictly increasing or strictly decreasing.

\subsubsection*{If \(f\) is increasing}
Assume \(f(x_0)>x_0\) for some \(x_0\). Then by monotonicity,
\[
x_0<f(x_0)<f^2(x_0)<\cdots<f^k(x_0)=x_0,
\]
impossible. Similarly \(f(x_0)<x_0\) is impossible. Thus \(f(x)=x\) for all \(x\), i.e. \(f=\mathrm{id}\).

But then
\[
S(x)=kx,
\]
which is not constant on \([0,1]\). So the increasing case gives no solution.

\subsubsection*{Therefore \(f\) is decreasing}
Then \(f^2\) is increasing, and some positive iterate of \(f^2\) is the identity, namely
\[
(f^2)^k=f^{2k}=\mathrm{id}.
\]
By the increasing-case argument, this forces
\[
f^2=\mathrm{id}.
\]
So \(f\) is a decreasing involution.

Also, since \(f^k=\mathrm{id}\), \(k\) must be even: if \(k\) were odd, then \(f^k\) would be decreasing, but the identity is increasing.

Thus write
\[
k=2m.
\]

\subsection*{3. Use the sum condition}

Because \(f^2=\mathrm{id}\), the iterates alternate:
\[
f_{2r}(x)=x,\qquad f_{2r+1}(x)=f(x).
\]
Hence
\[
S(x)=m\,x+m\,f(x)=m(x+f(x)).
\]
Since \(S(x)\) is constant, \(x+f(x)\) must be constant:
\[
x+f(x)=c \qquad \forall x\in[0,1].
\]
Therefore
\[
f(x)=c-x.
\]

Now \(f([0,1])\subseteq [0,1]\). From \(x=0\) we get \(c\in[0,1]\), and from \(x=1\) we get \(c-1\in[0,1]\), i.e. \(c\in[1,2]\). Hence
\[
c=1.
\]
So necessarily
\[
f(x)=1-x.
\]
Then
\[
S(x)=m(x+1-x)=m=\frac{k}{2}.
\]
Thus
\[
a=\frac{k}{2}.
\]

\subsection*{4. Existence}

For every even \(k\), the function
\[
f(x)=1-x
\]
is continuous from \([0,1]\) to \([0,1]\), and its iterates alternate:
\[
x,\ 1-x,\ x,\ 1-x,\dots
\]
So for \(k=2m\),
\[
\sum_{i=0}^{k-1} f_i(x)=m(x+1-x)=m=\frac{k}{2}.
\]
Hence such a function exists.


\medskip\hrule\medskip


\textbf{Final answer.}
\[
\boxed{\text{Such a continuous }f:[0,1]\to[0,1]\text{ exists exactly when }k\text{ is even and }a=\frac{k}{2}.}
\]
In that case one may take
\[
\boxed{f(x)=1-x.}
\]
\newpage
\end{document}